\section{Introduction}
\label{sec::intro}
Single-cell RNA sequencing (scRNA-seq) has emerged as a fundamental technology for profiling gene expression in bioinformatics, enabling high-resolution profiling of gene expression across individual cells. Based on the scRNA-seq data, cell clustering plays a central role by organizing cells into biologically meaningful groups based on their gene expression patterns, thereby revealing the intrinsic heterogeneity of cell populations and facilitating biological interpretation of their functions and interactions~\cite{kiselev2019challenges,peng2020single}. Accordingly, clustering scRNA-seq data is particularly valuable for a range of downstream tasks, including marker gene annotation~\cite{dai2022accurate} and cell type identification~\cite{hu2020iterative}.

For effectively scRNA-seq cell clustering, numerous methods have been proposed to achieve this goal. Early approaches primarily relied on classical clustering algorithms, such as $k$-means~\cite{hicks2021mbkmeans}, hierarchical clustering~\cite{johnson1967hierarchical}, and density-based methods~\cite{petegrosso2020machine}, which are computationally efficient but are inherently limited in modeling the complex nonlinear characteristics of scRNA-seq data. With the rapid advances in deep learning, recent studies have shifted toward learning-based clustering frameworks that leverage graph neural networks (GNNs)~\cite{ju2025survey,fan2026cmgl} to extract informative representations from scRNA-seq data. By representing cells as nodes and their interactions as edges in a graph, GNN-based approaches jointly exploit gene expression features and structure characteristics to learn context-aware cell embeddings, enabling a more comprehensive cell clustering~\cite{wang2021scgnn,gan2022deep,xu2024sccdcg,xu2025scsiameseclu,zhang2026evidence}.


\begin{figure}[t]
    \centering
    \includegraphics[width=0.9\linewidth]{figs/Meuro_human_Pancreas_cell_similarity_kde_black.png} 
    \caption{Similarity distributions of cell embeddings obtained by \method{} with and without structural information on the scHuman pancreas cells dataset.}
    \label{fig:similarity_distribution}
\end{figure}

Despite the effectiveness of these methods, we argue that two major challenges remain insufficiently addressed: \circlefir{} \textbf{\textit{Ignoring the Sparsity and Noise inherent within the scRNA-seq Data.}} Arising from extensive dropout effects and measurement noise, scRNA-seq data present a high-dimensional expression matrix~\cite {ou2022matrix,wu2022network}, which substantially complicates the reliable construction of cell relationships and frequently leads to graphs that are highly sparse and noisy. \circlesec{} \textbf{\textit{Limited Exploitation of Complex Intercellular Structural Information.}} While existing GNN-based approaches primarily focus on local neighborhood aggregation, they often overlook richer structural cues embedded within the cell graph. Features such as shortest paths and node distances are rarely leveraged, resulting in an insufficient utilization of the graph’s inherent information. As shown in Figure~\ref{fig:similarity_distribution}, cell similarities derived solely from gene expression features tend to be highly homogeneous and difficult to distinguish, whereas incorporating structural information such as positional relationships and shortest-path distances enables a more discriminative characterization of intercellular relationships for cell clustering.   

Having realized the above challenges with existing methods, in this paper, we propose a novel deep Siamese \underline{\textbf{G}}raph \underline{\textbf{T}}ransformer \underline{\textbf{N}}etwork for single-cell RNA sequencing (\textbf{\method{}}), which effectively exploits complementary graph information derived from gene expression data to learn discriminative and robust cell embeddings. Specifically, given the scRNA-seq data, we formulate it into a graph and construct two augmented graph views to serve as dual branches to capture and refine the intercellular information. Then, for each branch, we introduce a graph transformer network that incorporates shortest-path information and node-wise distances to capture the complementary graph structural dependencies, enabling a more expressive modeling of complex intercellular relationships. Based on the encoded information, we utilize an optimal transport-guided self-supervised clustering method to align cluster distributions for the clustering learning process.   

To summarize, the contributions of this paper are:
\begin{itemize}[leftmargin=*] 
% \vspace{-0.5pt} \setlength{\itemsep}{0pt} \setlength{\parsep}{-1pt} \setlength{\parskip}{0pt}

\item[\circlefir{}]  
\textbf{\textit{New Perspective:}} We study scRNA-seq data clustering by jointly considering gene expression profiles and cell graph structures, which effectively exploits complementary intercellular information beyond feature-level similarities. 

\item[\circlesec{}]  \textbf{\textit{Novel Methodology:}} We propose a novel framework termed \method{}, which introduces a Siamese graph transformer network to incorporate richer graph structural information, including shortest paths and node-wise distances for the cell clustering task.

\item[\circlethi{}] \textbf{\textit{Extensive Experiment:}} We conduct extensive experiments on multiple benchmark scRNA-seq datasets to evaluate the effectiveness of our \method{}. The results further show the superior clustering performance of our framework.
\end{itemize} 


% such as shortest paths and node distances, leading to insufficient prior information extracted from the graph..



% graph neural networks (GNNs) have gained increasing attention benefiting from the interaction structures 


% Thanks to the powerful ability of deep learning, recent works have increasingly adopted the deep learning framework for scRNA-seq data clustering. Benefiting from the power- ful utilization of graph information, graph neural networks (GNNs) have been applied to analyze scRNA-seq data by mod- eling cells as nodes and their interactions as edges, effectively capturing both gene expressions and cell graphs to comprehen- sively represent cellular heterogeneit Typically, graph neural network–based methods 

% naturally lend themselves to cell clustering downstream


% Spatial resolved transcriptomics techniques (SRT) measures gene expression while preserving the spatial location of cells within tissues, facilitating the investigation of previously challenging or impossible biological questions, such as the inherent heterogeneity within classical cell types, dynamic changes in cellular processes, reliable cell-cell communications, and the biological pathways underlying cellular differentiation~\cite{williams2022introduction}. Existing SRT data are mainly generated by two distinct protocols: barcode-based methods, such as 10x Visium~\cite{ji2020multimodal}, SLIDE-seq~\cite{rodriques2019slide}, and Stereo-seq~\cite{chen2022spatiotemporal}; and image-based methods, such as MERFISH~\cite{moffitt2018molecular}, osmFISH~\cite{codeluppi2018spatial}, and seqFISH~\cite{lubeck2014single,shah2016situ}. Unlike single-cell resolution sequencing such as scRNA-seq, the resolution of these SRT platforms is typically at the spot level (ranging from a few to tens of cells). Despite this limitation, SRT has been widely applied in biological research, generating vast amounts of data, which urgently require computational methods for data mining.

% Among the key analytical tasks in SRT data, spatial domain identification, which aims to reveal the spatial distribution of different cellular populations or subpopulations, and the architectural features of pathological tissues, has garnered significant attention. It is also considered a foundation for various downstream tasks, such as cell type annotation, differential gene expression analysis, and trajectory inference. Early approaches primarily adapted traditional unsupervised algorithms, including PCA~\cite{abdi2010principal}, Seurat~\cite{satija2015spatial}, with clustering methods, such as K-means~\cite{likas2003global} and mclust~\cite{fraley2002model}. However, these methods treat cells as independent data points, disregarding the rich biological prior knowledge and cellular dependencies, leading to the loss of valuable information. In recent years, Graph Neural Networks (GNNs) have been widely employed to capture the intercellular dependency network, which has advanced the development of spatial domain identification algorithms. SpaGCN~\cite{hu2021spagcn} pioneers domain identification with graph convolutional network, enabling spatially variable genes analysis. DeepST~\cite{xu2022deepst} computes the graph adjacency matrix based on spatial coordinates by $k$-nearest neighbours and incorporates pathological image information to enhance spatial domain detection performance. STAGATE~\cite{dong2022deciphering} adopts a graph attention mechanism to adaptively learn the similarity of neighboring spots. GraphST~\cite{long2023spatially} combines graph neural networks with self-supervised contrastive learning, enabling clustering, integration, and deconvolution of SRT data, while stHGC~\cite{wang2024sthgc} builds on a hybrid neighbor graph constructed by integrating different similarity metrics.

% The aforementioned graph-based methods have demonstrated superiority in spatial domain identification tasks, however, there are still several intrinsic limitations to be addressed. \emph{\textbf{In particular}}, most existing methods \textbf{simplify the complex intercellular dependencies by abstracting them into $k$-nearest neighbor graph}, which results in the loss of amount of valuable information. Some of the methods construct graph structures based on spatial coordinates~\cite{xu2022deepst,long2023spatially,oh2025global}, overlooking the fact that spots of the same type or subtype are not always spatially adjacent, while others attempts to alleviate information loss by incorporating modal information beyond spatial proximity~\cite{dong2022deciphering,hu2021spagcn,wang2024sthgc}, which still fail to capture the hierarchical and structural information of cellular dependency networks. \emph{\textbf{Moreover}}, %the spatial domain identification task can be divided into two modules: representation learning and cell/spot clustering. Most of the previous methods treat these two modules as separate components \textbf{without adequately addressing the need to effectively couple the representation learning process with clustering}~\cite{dong2022deciphering, xu2022deepst, long2023spatially, wang2024sthgc}. This lack of integration prevents these methods from constraining the model to simultaneously learn representations that are conducive to clustering while also generating accurate and coherent cluster assignments, leading to suboptimal clustering performance. \emph{\textbf{Finally}}, 
% \textbf{spots located at the domain boundary are frequently misassigned to incorrect clusters}~\cite{dong2022deciphering,xu2022deepst,long2023spatially}. These boundary-located spots, characterized by higher levels of uncertainty, could benefit from the integration of biological prior knowledge, which may provide valuable contextual information. By leveraging such domain-specific insights, it is possible to improve the assignment of these spots to the correct clusters, thereby enhancing the overall accuracy and biological relevance of the clustering process.

% To address these issues, we propose the \underline{Bio}logy-informed \underline{H}yper\underline{G}raph \underline{N}eural \underline{N}etwork for domain identification in spatially resolved transcriptomics data (\method{}). We begin by constructing hypergraphs corresponding to the information from the spatial and feature views to represent the hierarchical cellular dependencies. Subsequently, two hypergraph neural networks, with non-shared parameters, are used as encoders to learn low-dimensional representations. We fuse multi-view embeddings through a cross-attention mechanism and apply a global self-correlation strategy for representation enhancement. To constrain the model to learn clustering-friendly representations, we adopt the triplet self-supervised loss, while the hypergraph smoothing loss is used to alleviate the uncertainty in boundary-located spots' cluster assignments. Finally, to handle the high sparsity and over-dispersion of SRT data, we utilize the zero-inflated negative binomial (ZINB) decoder to reconstruct the gene expression matrix. Extensive qualitative and quantitative experiments, as well as downstream analyses, are conducted on benchmark datasets from different species and organs, demonstrating the effectiveness and superiority of our proposed \method{}.